% No preamble. Prefix jsi-. Ordered scalar limit and exact holonomy obstruction.
\section{The scalar infrared limit and a forced holonomy cancellation}
\label{jsi-section}

There are two different assertions in this section. A common positive
physical mass gives a well-defined thermodynamic scalar reference,
whose subsequent massless limit is \(\alpha_g\log K\). The original
separately massless globally based determinants need not have that
order-independent analytic limit. An exactly periodic flat connection
gives a concrete obstruction and, through the actual finite determinant
identity, forces a compensating response in the physical factor.

\subsection{The mass convention and an exact base correction}
\label{jsi-mass-section}

Fix \(K=\ell^k\) and final torus side \(N\) in four dimensions.
Current fine sites have physical spacing \(K^{-1}\), with side \(KN\);
coarse sites have spacing one and side \(N\).
All color coordinates are trace orthonormal.
Let \(G_{\rm all}(U)\) be the covariant scalar Laplacian on all sites,
and let \(G_b(U)\) be its principal minor after deleting the color
coordinates at one common physical base \(o\). Deleting a site does
not change the remaining diagonal entries.

A common physical mass \(\delta>0\) means
\begin{equation}
 H_{f,\delta}=K^2G_{{\rm all},f}+\delta^2I,\qquad
 H_{c,\delta}=G_{{\rm all},c}+\delta^2I.
 \label{jsi-mass}
\end{equation}
Their principal minors are \(H_{f,\delta}^b,H_{c,\delta}^b\).
With \(\mathcal N_0F=F(W)-F(1)\), define
\begin{equation}
 S_{\delta,N}^b=\frac12\mathcal N_0
       \bigl(\log\det H_{f,\delta}^b-\log\det H_{c,\delta}^b\bigr).
 \label{jsi-scalar-ratio}
\end{equation}
The fine counting mass is \(\delta/K\), not \(\delta\).
At zero mass multiplication of the fine operator by \(K^2\)
changes only a background-independent normalization.
Equation \eqref{jsi-scalar-ratio} is a specified scalar comparison
family. It is not automatically an exact mass deformation of the
entire physical-quotient determinant identity.

\begin{lemma}[Positive-mass base defect]
\label{jsi-base-lemma}
Put \(H=\epsilon^{-2}G_{\rm all}+\delta^2I\), \(R=H^{-1}\),
where \(\epsilon=K^{-1}\) or \(1\). Then
\begin{equation}
       \det H^b=\det H\,\det R_{oo}.
 \label{jsi-jacobi}
\end{equation}
At fixed \(\delta>0,K\), the normalized logarithm of the fixed color
block \(R_{oo}\) has a complex product neighborhood and first-two
bounded-direction derivative bounds independent of \(N\).
Its first-two source kernels are exponentially summable in their
distances from \(o\), including all fixed polynomial moments.
No uniformity as \(\delta\downarrow0\) is asserted.
\end{lemma}

\begin{proof}
The determinant identity is the block Jacobi identity. At real
backgrounds,
\[
 \delta^2I\le H\le(16\epsilon^{-2}+\delta^2)I,
 \qquad
 (16\epsilon^{-2}+\delta^2)^{-1}I
                      \le R_{oo}\le\delta^{-2}I.
\]
The nearest-neighbor coefficients have volume-independent operator
bounds on a small complex product ball. The resolvent identity,
and then the lower bound for \(R_{oo}\), give a smaller complex
ball on which both inverses remain bounded uniformly in volume.
The logarithm branch continued from zero is bounded there after
normalization. Its first two derivatives also follow directly from
\begin{align}
 DR[v]&=-RH'[v]R,\notag\\
 D^2R[v,w]&=RH'[v]RH'[w]R+RH'[w]RH'[v]R-RH''[v,w]R.
 \label{jsi-resolvent-derivatives}
\end{align}
Differentiating \(\log\det R_{oo}\) uses only its fixed color trace.

For exponential bounds, conjugate by a site weight whose ratios
across nearest-neighbor edges differ from one by \(O(\kappa)\).
Choose \(\kappa>0\), depending on \(\delta,\epsilon\), so the
conjugation perturbation has norm at most \(\delta^2/4\).
The weighted inverse is then bounded by \(2/\delta^2\).
Spending a smaller exponent gives absolute row and column Schur
bounds. The derivatives \(H',H''\) are local. Every source chain
in \eqref{jsi-resolvent-derivatives}, and in the logarithmic
derivatives, begins and ends at \(o\). Its exponential weights
therefore control all source distances from \(o\), with a further
exponent reserve for any fixed polynomial moment.
\end{proof}

Writing \(S_{\delta,N}^{\rm all}\) for the unbased half-ratio gives
the exact comparison
\begin{equation}
 S_{\delta,N}^b=S_{\delta,N}^{\rm all}
 +\frac12\mathcal N_0\bigl\{
       \log\det(R_f)_{oo}-\log\det(R_c)_{oo}\bigr\}.
 \label{jsi-defect}
\end{equation}
The lemma has two useful consequences. A base receding from fixed
localized sources changes their first-two coefficients
exponentially little. Alternatively, average the based functional
over the \(N^4\) simultaneous coarse translations of the fine and
coarse base. Reanchoring its Hessian at one output cell bounds the
anchored second-moment norm of its difference from the unbased
functional by
\begin{equation}
                         C_{\delta,K}N^{-4}.
 \label{jsi-average-defect}
\end{equation}
Indeed summing the translated rooted defect kernel converts that
row into its full source sum, divided by \(N^4\); its displacement
moments are bounded by the lemma. This averaging is an estimate
for the limit, not a claim that different finite base choices
give identical functionals.

After composition with the actual fixed-\(K\) stationary branch,
the same conclusions use its previously proved localized source
kernels. Under the Notebook 33 inputs the constants are uniform
in old depth \(M\), still depending on \(K,\delta\).

\subsection{The ordered localized Maxwell limit}
\label{jsi-ordered-section}

At positive mass, the unbased diagonal logarithm has exponentially
localized Taylor kernels, by the same resolvent argument applied
to its logarithm integral. Its color trace is gauge invariant.
Its Hessian moments therefore define its affine curvature
coefficient absolutely. This gives a way to take the limit without
postulating an affine magnetic field on a torus: first pass to
infinite kernels using compactly supported directions, then
evaluate polynomial directions by their exponential moments.
Symmetric short lifts of periodic kernels give the same result.
At fixed \(K,\delta\), unbased winding contributions vanish
exponentially in \(N\), and \eqref{jsi-average-defect} removes the
averaged base contribution.

For an actual stationary insertion, use fixed-\(K\) structural
compactness followed by the polynomial-jet identification of
Notebook 33. Every such limit has the same scalar expression
below. This statement uses the already established source kernels;
it does not infer locality from a finite determinant identity.

Normalize the Maxwell functional by
\(\beta(F)=2[t^2]\) of its localized constant-curvature density.
For \(\mu\ne\nu\) and \(\tau(H^2)=1\), the current and coarse
Cartan probes are
\[
 U_\mu(t;x)=e^{itx_\nu H/K^2},\qquad
 W_\mu(t;y)=e^{ity_\nu H}.
\]
The unpinned scalar heat coefficient already established is
\begin{equation}
 \begin{split}
 h(s)&=\frac{C_H}{3}s^2e^{-8s}I_0(2s)^2I_1(2s)^2,\\
 h(s)&=O(s^4)\quad(s\downarrow0),\\
 h(s)&=\alpha_g+O(s^{-1})\quad(s\to\infty),
 \qquad \alpha_g=\frac{C_H}{48\pi^2}.
 \end{split}
 \label{jsi-heat}
\end{equation}
The \(K^4\) fine sites per final block cancel the curvature
factor \(K^{-4}\). In common physical heat time \(u\), the fine
and coarse coefficients are consequently \(h(K^2u)\) and \(h(u)\).

\begin{proposition}[Ordered scalar reference limit]
\label{jsi-ordered-theorem}
Take the thermodynamic localized-kernel limit first, at fixed
\(\delta>0\) and fixed \(K\). The resulting scalar coefficient is
\begin{equation}
 \beta_\delta(K)=\frac12\int_0^\infty
           e^{-\delta^2u}\{h(K^2u)-h(u)\}\frac{du}{u}.
 \label{jsi-beta-delta}
\end{equation}
It satisfies
\begin{equation}
             \lim_{\delta\downarrow0}\beta_\delta(K)
                         =\alpha_g\log K.
 \label{jsi-beta-limit}
\end{equation}
For \(K\ge1\) and \(0<\delta\le1/2\), the explicit reference
also obeys
\begin{equation}
 |\beta_\delta(K)-\alpha_g\log K|
       \le C\delta^2\{1+\log(1/\delta)\},
 \label{jsi-mass-rate}
\end{equation}
where \(C\) is independent of \(K\).
\end{proposition}

\begin{proof}
The localized logarithm heat formula, the preceding normalization,
and \eqref{jsi-heat} give \eqref{jsi-beta-delta}.
Without the mass factor, the paired integrand is absolutely
integrable: use its small-\(s\) bound at zero and cancellation
of \(\alpha_g\) at infinity. On \(\varepsilon<u<T\), its
unhalved integral equals
\[
 \int_T^{K^2T}h(v)\frac{dv}{v}
       -\int_\varepsilon^{K^2\varepsilon}h(v)\frac{dv}{v}.
\]
The limits are \(2\alpha_g\log K\) and zero, respectively.
Dominated convergence proves \eqref{jsi-beta-limit}.

For the rate, \(h\) is bounded on the whole positive axis.
On \(0<u<1\), use \(1-e^{-\delta^2u}\le\delta^2u\);
the resulting error is \(O(\delta^2)\), uniformly in \(K\).
For \(u\ge1\), \eqref{jsi-heat} gives
\(|h(K^2u)-h(u)|\le2C/u\) for every \(K\ge1\).
Split the remaining integral at \(u=\delta^{-2}\).
The first part is \(O(\delta^2\log(1/\delta))\), the
second \(O(\delta^2)\).
\end{proof}

The uniformity in \eqref{jsi-mass-rate} belongs to the explicit
reference after the thermodynamic limit. It does not interchange
mass removal with that limit for the original based determinants.

\subsection{An exact massless anomaly on a periodic flat connection}
\label{jsi-anomaly-section}

There is a direct obstruction to such an interchange, already
at zero momentum. Consider one charged scalar \(q\) on a periodic
\(n^4\) lattice, with constant phase \(e^{itqA}\) on direction
\(\mu\). This is a legitimate periodic flat connection for every
\(t\); no continuously varying magnetic flux is involved.
Its unbased eigenvalues are
\begin{equation}
 \lambda_p(t)=2\sum_\nu
       \{1-\cos(p_\nu+tqA\,\delta_{\mu\nu})\},
       \qquad p\in(2\pi/n)\mathbb Z_n^4.
 \label{jsi-spectrum}
\end{equation}
Put \(V=n^4\), and delete one scalar site to obtain \(G_n^b(t)\).
Translation invariance of the resolvent diagonal and
\eqref{jsi-jacobi} imply
\begin{equation}
 \det G_n^b(t)=\frac1V\prod_{p\ne0}\lambda_p(t)
       \left\{1+\lambda_0(t)
                    \sum_{p\ne0}\lambda_p(t)^{-1}\right\}.
 \label{jsi-cofactor}
\end{equation}
Initially this is an identity away from zero eigenvalues;
both sides continue analytically through \(t=0\).

\begin{proposition}[Noncommuting mass and thermodynamic derivatives]
\label{jsi-anomaly-theorem}
Let
\begin{equation}
 g_4=\int_{[-\pi,\pi]^4}\frac{dp}{(2\pi)^4}
             \frac1{2\sum_\nu(1-\cos p_\nu)}>0.
 \label{jsi-green}
\end{equation}
Then
\begin{equation}
 \lim_{n\to\infty}n^{-4}[t^2]\log\det G_n^b(t)
                           =q^2A^2g_4.
 \label{jsi-anomaly}
\end{equation}
At every fixed positive scalar mass, the corresponding
thermodynamic coefficient is zero, and remains zero when
the mass is subsequently removed.
\end{proposition}

\begin{proof}
Since \(\lambda_0(t)=q^2A^2t^2+O(t^4)\),
\eqref{jsi-cofactor} gives, writing \(\lambda_p=\lambda_p(0)\),
\begin{equation}
 [t^2]\log\det G_n^b(t)
   =q^2A^2\sum_{p\ne0}
    \left\{\frac{\cos p_\mu}{\lambda_p}
       -\frac{2\sin^2p_\mu}{\lambda_p^2}
       +\frac1{\lambda_p}\right\}.
 \label{jsi-cofactor-jet}
\end{equation}
The first two summands equal
\(\partial_{p_\mu}(\sin p_\mu/\lambda_p)\).
Their singularity is \(O(|p|^{-2})\), integrable in dimension
four. Removing a radius-\(\varepsilon\) ball before integration
by parts leaves boundary flux \(O(\varepsilon^2)\), so their
Brillouin-zone integral is zero. Their grid sums converge:
\begin{equation}
 n^{-4}\sum_{0<|p|\le\varepsilon}|p|^{-2}
                       \le C\varepsilon^2+O(n^{-2}).
 \label{jsi-grid-bound}
\end{equation}
Outside that ball ordinary Riemann convergence applies.
The same argument for \(\lambda_p^{-1}\) gives \(g_4\),
proving \eqref{jsi-anomaly}.

At fixed positive mass \(m\), the unbased coefficient per volume
instead tends to the integral of
\(\partial_{p_\mu}\{\sin p_\mu/(\lambda_p+m^2)\}\),
which is zero by periodicity. Lemma~\ref{jsi-base-lemma}
makes the based correction \(O(1)\) in these derivatives,
so division by \(n^4\) removes it. Subsequent \(m\downarrow0\)
still gives zero.
\end{proof}

The deleted block has fixed rank, but its massless derivative
contribution is extensive. Equation \eqref{jsi-cofactor} identifies
the soft global mode responsible for the failure. Summing
adjoint charges replaces \(q^2\) by \(C_H\); for a vector of
commuting constant angles the answer is \(C_H|A|^2g_4\),
with mixed terms zero by coordinate reflections.

\subsection{The actual branch and its forced physical compensation}
\label{jsi-actual-section}

Now use the actual final constant links and current representative
\begin{equation}
 W_\mu(t)=e^{itA_\mu H},\qquad
 U_\mu(t)=e^{itA_\mu H/K}.
 \label{jsi-flat-branch}
\end{equation}
Each ordinary blocking stage preserves this commuting connection:
all paths with the same displacement have the same transporter,
so the selected mean is evaluated on identical entries.
The same flat lift exists through every old refinement \(M\).
All plaquettes are identity. Since the real Wilson action is
nonnegative, the associated classical minimum values are zero.
After the actual local root gauge, uniqueness of the small
minimum germ identifies this with the actual stationary branch.
This is a fixed finite-germ argument.

Every sitewise adjoint gauge preserves a scalar zero at the
global base. Thus the globally based scalar determinants are
unchanged by that root gauge and may be computed in
\eqref{jsi-flat-branch}.
Fine current volume is \((KN)^4\) and fine angle is \(tA/K\);
coarse volume is \(N^4\) and angle is \(tA\).
Proposition~\ref{jsi-anomaly-theorem} therefore gives
\begin{equation}
 \lim_{N\to\infty}N^{-4}[t^2]
         (\log\det G_f-\log\det G_c)
             =C_H|A|^2(K^2-1)g_4.
 \label{jsi-orbit-anomaly}
\end{equation}
The generators on this branch are the same for every old \(M\).
This is a flat-holonomy, zero-momentum response, not a Maxwell
coefficient.

\begin{proposition}[Forced compensation in the physical factor]
\label{jsi-compensation-theorem}
Fix \(K\) before taking the thermodynamic limit.
For the actual leading correction and physical quotient,
\begin{align}
 \lim_{N\to\infty}N^{-4}[t^2]D_N(W(t))&=0,
       \label{jsi-total-zero}\\
 \lim_{N\to\infty}N^{-4}[t^2]\log\det Q_\perp(W(t))
                  &=C_H|A|^2(K^2-1)g_4.
       \label{jsi-compensation}
\end{align}
The same conclusion allows varying old depth under the
previously established uniform-old-depth, fixed-\(K\) local
kernel bounds.
\end{proposition}

\begin{proof}
The actual correction \(D_N\) has translation-covariant Hessian
kernels with common exponential row and source bounds at fixed
\(K\). Every structural limit has invariant local densities and
zero linear color jet. Its quadratic Ward identity consequently
annihilates lattice gradients. Exponential moments extend this
identity to linear scalar gauge parameters by cutoff approximation.
Their gradients are constant link fields, so the infinite
quadratic symbol vanishes at momentum zero.

The finite normalized coefficient on \eqref{jsi-flat-branch} is
one-half the contracted Hessian row sum. Uniform exponential
tails make its distant part uniformly small. Every fixed
central sum converges along a structural subsequence, and its
full limiting row sum is zero by the preceding Ward argument.
All subsequential scalar limits are therefore zero, proving
\eqref{jsi-total-zero}. The uniform-old-\(M\) versions of these
same bounds permit \(M\) to vary. No uniform-in-\(K\) locality
is used.

At finite \(N\) the exact physical-quotient identity is
\begin{equation}
 2D_N=\mathcal N_0
       (\log\det G_f-\log\det G_c-\log\det Q_\perp).
 \label{jsi-actual-identity}
\end{equation}
Subtract \eqref{jsi-total-zero} from
\eqref{jsi-orbit-anomaly} using the same branch, base and
regulator. This yields \eqref{jsi-compensation}, including
its sign and normalization.
\end{proof}

No claim that \(D_N(W(t))\) is identically zero at finite \(N\)
is needed: a periodic constant connection may have nontrivial
holonomy. Nevertheless, its thermodynamic constant-field
response must vanish by locality and the Ward identity.
Thus the term \(-\frac12\log\det Q_\perp\) cancels the scalar
half-ratio exactly at zero momentum.

\subsection{What the ordered limit does not identify}
\label{jsi-scope}

The positive result is
\[
 \lim_{\delta\downarrow0}\lim_{N\to\infty}
       \{\hbox{localized regulated scalar Maxwell coefficient}\}
                         =\alpha_g\log K.
\]
The original separately massless globally based family does not
have an order-independent analytic local limit, as the periodic
test \eqref{jsi-orbit-anomaly} proves.
In particular a \(K\)-uniform \(O(1)\) constant-potential response
for the separate physical determinant is false, despite its
flat lower bound \(cK^{-2}\).

The anomaly alone neither computes nor disproves equality of
Maxwell coefficients after a suitable coupled cancellation.
A valid joint vector/scalar regulator must preserve the
rank-fixed but derivative-large base factors before its
infrared limit. The exact finite quotient identity does not
provide that regulated identity by itself. No continuum
measure, nontriviality, or mass-gap conclusion follows here.
